(1) 由 \(P=K\rho^2\) 得 \(2K\rho\,\rho'=-Gm\rho/r^2\)，在 \(\rho>0\) 区域约去密度，得 \(2K\rho'=-Gm/r^2\)。对 \(r^2\rho'\) 求导并用 \(m'=4\pi r^2\rho\)，得到
\[
\frac1{r^2}(r^2\rho')'+\frac{2\pi G}{K}\rho=0.
\]
令 \(a=\sqrt{K/(2\pi G)}\)、\(\xi=r/a\)，正则解为
\(\rho=\rho_c\sin\xi/\xi\)；另一个 \(\cos\xi/\xi\) 分支中心奇异而被排除。第一零点 \(\xi=\pi\) 给 \(R=\pi a\)。

(2) 直接积分 \(\rho_c\sin(r/a)/(r/a)\)：
\[
m(r)=4\pi a^3\rho_c(\sin\xi-\xi\cos\xi),\qquad M=4\pi^2a^3\rho_c .
\]
当 \(\xi\to0\)，\(\sin\xi-\xi\cos\xi=\xi^3/3+O(\xi^5)\)，故 \(m\sim4\pi\rho_cr^3/3\) 保证中心引力有限。表面引力大小 \(GM/R^2=4G a\rho_c\)，且 \(P(R)=K\rho(R)^2=0\)，与无外压边界相容。

(3) 先算压强体积分：
\[
\Pi=4\pi\int_0^R P r^2dr
=4\pi K\rho_c^2a^3\int_0^\pi\sin^2\xi\,d\xi
=2\pi^2K\rho_c^2a^3.
\]
从 \(\int 4\pi r^3P'(r)dr\) 分部积分并用 \(P(R)=0\)，得到 \(W=-3\Pi=-6\pi^2K\rho_c^2a^3\)。另由上问 \(M=4\pi^2a^3\rho_c,R=\pi a,K=2\pi Ga^2\) 直接化简，验证
\(\boxed{W=-3GM^2/(4R)}\)，这是独立的引力能量尺度核验。

(4) 对 \(\gamma=2\)，流体内部能为 \(U=\int P/(\gamma-1)dV=\Pi\)，于是静态总能量 \(E=U+W=-2\Pi=-4\pi^2K\rho_c^2a^3<0\)。固定 \(K,G\) 时 \(a,R\) 与 \(\rho_c\) 无关，\(M\propto\rho_c\)、\(E\propto-\rho_c^2\)，这是 \(n=1\) 模型的特殊尺度性而非所有恒星的普遍性质。中心密度过高会破坏弱场牛顿近似，若 \(GM/(Rc^2)\) 不再远小于1，则必须引入广义相对论结构方程。
